\section{前馈式统一几何预测与方法趋势}

\subsection{从可微优化到统一前馈推断}

在可微 SfM 之后，VGG 团队研究主线继续向“统一化”与“前馈化”推进。\textit{VGGT: Visual Geometry Grounded Transformer} 提出以单个前馈网络直接预测相机参数、深度图、点图和三维点轨迹等关键几何属性，试图将传统多阶段系统进一步压缩为统一模型\cite{wang2025vggt}。这表明研究问题的表述已经从“如何提升某一几何模块”转向“如何以统一表示组织多个几何任务”。

VGGT 的意义不只在于预测速度更快，更在于它反映出多视图几何重建研究的一个重要趋势：学习模型开始承担传统后端优化原本负责的更大一部分工作。相比于需要测试时迭代求解的系统，前馈式方法在部署、复用和跨任务迁移上具有更高潜力\cite{wang2025vggt}。这也是 VGGT 被视为 VGG 团队研究主线阶段性代表作的重要原因。

\subsection{与相关方法的横向比较}

为了更好理解 VGGT 的定位，可以将其与 DUSt3R 等近期方法进行对照。DUSt3R 也强调从图像中直接恢复三维几何，并在一定程度上简化传统结构恢复流程\cite{wang2024dust3r}。但与 VGG 团队这一主线相比，VGGT 更强调统一的几何属性输出，以及与前述几何先验、轨迹学习和可微 SfM 研究的连续性。

此外，\textit{Splatter Image: Ultra-Fast Single-View 3D Reconstruction} 虽然聚焦单图重建，但它同样体现出 VGG 团队对“前馈式三维表示”的持续探索\cite{szymanowicz2024splatter}。这说明 VGG 团队并非只在传统多视图任务上推进学习化，也在尝试寻找更统一、更高效的三维表示方式。

\subsection{向动态几何与统一表示扩展}

这一趋势还继续延伸到动态几何场景。\textit{Dynamic Point Maps: A Versatile Representation for Dynamic 3D Reconstruction} 说明，当场景随时间变化时，统一点图表示仍可作为多任务几何建模的重要中间层\cite{sucar2025dpm}。尽管这类工作已部分超出本文“多视图几何重建”的核心范围，但它们揭示了一个清晰方向：未来的学习方法可能不再局限于重建某个静态结果，而是形成可支持深度、相机、点云、轨迹乃至动态场景理解的统一几何表征。

\subsection{本阶段的综合评价}

从 VGG 团队近五年的方法演进来看，前馈式统一几何预测并不是凭空出现的，而是建立在前述几何先验注入、点轨迹学习与可微 SfM 逐步成熟的基础之上。也正因为如此，VGGT 等方法更适合被理解为一条研究主线的阶段性总结，而不是孤立的单篇突破。对本文主题而言，本阶段最值得强调的不是单个模型的性能，而是“多视图几何重建正在从优化驱动转向表示驱动与模型驱动”的方法论变化。
